\documentclass[10pt,a4paper]{amsart}
\usepackage[margin=19mm]{geometry}
\usepackage{amsmath,amssymb,mathtools}
\usepackage[scr=rsfs,cal=euler]{mathalfa}
\usepackage{xfrac}
\usepackage[unicode,hidelinks,bookmarks=false]{hyperref}
\newcommand{\ie}{i.e.\ }
\newcommand{\cf}{cf.\ }
\newcommand{\resp}{respectively\ }
\newcommand{\Subsection}{\subsection}
\providecommand{\nobreakdash}{\nobreak\hbox{-}}
\setlength{\emergencystretch}{2em}
\multlinegap0pt
\usepackage{bm}
\usepackage{float}
\usepackage{xcolor}
\usepackage{tikz}
\newtheorem{lemma}{Lemma}
\def\Ac{\ensuremath{\mathcal{A}}}
\newcommand{\E}{\mathbb{E}}
\newcommand{\St}{\mathcal{S}}
\def\evc{{c}}
\def\evmu{{\mu}}
\def\mB{{\bm{B}}}
\def\mP{{\bm{P}}}
\def\mTg{{\bm{T}}_\gamma}
\def\vc{{\bm{c}}}
\def\ve{{\bm{e}}}
\def\vmu{{\bm{\mu}}}
\def\vu{{\bm{u}}}
\def\vx{{\bm{x}}}
\def\vy{{\bm{y}}}
\title{Apprenticeship learning with prior beliefs using inverse optimization}
\author{Mauricio Junca, Esteban Leiva}
\date{Source: arXiv:2505.21639v2 (CC BY 4.0)}
\begin{document}
\maketitle

\noindent\emph{Source and licence.} Apprenticeship learning with prior beliefs using inverse optimization. Version arXiv:2505.21639v2 (CC BY 4.0), distributed by arXiv under the Creative Commons Attribution 4.0 International licence, http://creativecommons.org/licenses/by/4.0/, as recorded in the arXiv licence metadata for this exact version and shown on the abstract page https://arxiv.org/abs/2505.21639v2. Full licence notice: \texttt{licenses/arxiv-2505.21639-CC-BY-4.0.txt}.\par\smallskip

\noindent\textit{Editorial scope.} Counted unit: the article's lemma lemma:cu\_grad (source0.tex lines 412--415). The article's complete original proof of it is delivered with it (source0.tex lines 760--800, one \texttt{proof} environment of 41 lines, which the article places away from the statement). No mathematics is written, repaired or re-derived: the statement and the proof are the article's own lines, in the article's own notation, and no retained line is substituted or re-flowed.  No further source-local context is needed: every cross-reference of every retained line resolves inside the two delivered spans.  Nothing was omitted from any retained span, so the package carries no omission record.  No retained line contains a citation, so the wrapper carries no bibliography and no bibliography entry.  No retained line contains a figure environment, an image-inclusion command or drawing code, so no image is delivered and the package declares no asset dependency.  Editorial formatting only, in addition: an AMS \texttt{amsart} preamble that restates the article's own declarations and macros used by the retained lines, namely the environments \texttt{lemma} on separate counters, together with the standard packages those lines need.  The retained environments print in this document as Lemma 1 in order of appearance, whereas the article prints the same items with the numbering of its own full document.  The complete native source of the article as distributed in the arXiv e-print for this exact version is bundled unchanged as \texttt{sources/178878/source0.tex}.
\begin{lemma} \label{lemma:cu_grad}
    Gradient estimator $\Tilde{g}^{(\vc,\vu)}((\vc_t,\vu_t),\vmu_t)$ is a $(v^{(\vc,\vu)}, \Vert \cdot \Vert_2)$-bounded estimator, with 
    $$v^{(\vc,\vu)} =  64\alpha^2\cdot |\St||\Ac| + \frac{4(1+\gamma^2)}{(1-\gamma)^2} + 8.$$   
\end{lemma}

\begin{proof}[Proof of Lemma \ref{lemma:cu_grad}]
        The unbiasedness follows from the following observations
        \begin{align*}
            &\sum\limits_{s,a}\frac{|\St||\Ac|}{|\St||\Ac|}\cdot 2\alpha\left(\evc_t(s,a)\ve^{(s,a)}-\hat{\evc}(s,a)\ve^{(s,a)}\right)\\
            &+\sum\limits_{s_E,a_E} \evmu_{\pi_E}(s_E,a_E)\ve^{(s_E,a_E)} - \sum\limits_{s_t,a_t}\evmu_t(s_t,a_t)\ve^{(s_t,a_t)}\\
            &= 2\alpha(\vc_t - \hat{\vc}) + \vmu_{\pi_E} - \vmu_t
        \end{align*}
        and
        \begin{align*}
            \frac{1}{(1-\gamma)} &\sum\limits_{s_t',a_t,s_t} \evmu_t(s_t,a_t)P(s^\prime_t\:|\:s_t,a_t)(\ve^{(s_t)} - \gamma \ve^{(s_t^\prime)})\\
            &= \frac{1}{(1-\gamma)}\left(\sum\limits_{s_t,a_t}\evmu_t(s_t,a_t)\ve^{(s_t)} -\gamma \sum\limits_{s_t',a_t,s_t} \evmu_t(s_t,a_t)P(s^\prime_t\:|\:s_t,a_t) \ve^{(s_t^\prime)} \right)\\
            &= \frac{1}{(1-\gamma)}(\mB\vmu - \gamma\mP\vmu)\\
            &= \mTg\vmu.
        \end{align*}
        Thus, we obtain
        \begin{align*}
        \E\left[ \Tilde{g}^{(\vc,\vu)}((\vc,\vu),\vmu) \right] &=                                                 \begin{pmatrix}
                    2\alpha(\vc_t-\hat{\vc}) + \vmu_{\pi_E} - \vmu_t \\
                    \mTg\vmu - \mTg\vmu_{\pi_E}
                    \end{pmatrix}.
        \end{align*}
        To prove the bound on the second-moment, note that
        \begin{align*}
            &\E\left[\Big\Vert |\St||\Ac|\cdot 2\alpha\left(\evc_t(s,a)\ve^{(s,a)}-\hat{\evc}(s,a)\ve^{(s,a)}\right) + \ve^{(s_E,a_E)} - \ve^{(s_t,a_t)}\Big\Vert_2^2\right] \\
            &\leq 2\E\left[4|\St||\Ac|\alpha^2\Big\Vert\evc_t(s,a)\ve^{(s,a)}-\hat{\evc}(s,a)\ve^{(s,a)} \Big\Vert_2^2 + \Big\Vert \ve^{(s_E,a_E)} - \ve^{(s_t,a_t)}\Big\Vert_2^2\right]\\
            &\leq 2\E\left[4|\St||\Ac|\alpha^2\cdot4 + 2\right]\\
            &= 32|\St||\Ac|\alpha^2 + 4,
        \end{align*}
        and 
        \begin{align*}
            \E\left[\Big\Vert \frac{1}{(1-\gamma)}\left( \ve^{(s_t)} - \gamma \ve^{(s_t^\prime)} - (\ve^{(s_E)} - \gamma\ve^{(s_E^\prime)})\right) \Big\Vert_2^2\right] &\leq \E\left[\frac{2(1+\gamma^2)}{(1-\gamma)^2}\right]\\
            &= \frac{2(1+\gamma^2)}{(1-\gamma)^2}.
        \end{align*}
        Hence, we can provide the bound
        \begin{align*}
            \E\left[ \big\Vert \Tilde{g}^{(\vc,\vu)}((\vc,\vu),\vmu) \big\Vert^2_2 \right]
            &\stackrel{(i)}{\leq} 2\left[ 32|\St||\Ac|\alpha^2 + 4 + \frac{2(1+\gamma^2)}{(1-\gamma)^2}\right] \\
            &= 64\alpha^2\cdot |\St||\Ac| + \frac{4(1+\gamma^2)}{(1-\gamma)^2} + 8.
        \end{align*}
        where in $(i)$ we used $\Vert \vx+ \vy\Vert^2 \leq 2[\Vert \vx \Vert^2 + \Vert \vy \Vert^2]$.
     \end{proof}

\end{document}
